% ======================================================================
% RESULTS — six subsections mapped to Fig. 1–5.
% Structure and claim-gating only; every quantitative statement is a
% placeholder. No experimental results are invented.
% ======================================================================
\section*{Results}

% ----------------------------------------------------------------------
\subsection*{System design for embodied perception of subtle mechanical stimuli}
% -> Fig. 1
The tactile skin integrates a passive magnetic cilia layer, a soft elastomer/
adhesive interface, a distributed MLX90393 tri-axis magnetic sensing layer and
a hand-shaped flexible printed circuit (Fig.~1). Cilia are confined to the
palmar side---palm, finger pads and fingertips---so that the sensing surface
matches the contact surface of a grasping hand (Fig.~1a). Weak mechanical
inputs deflect the cilia, producing distinct three-axis magnetic signatures
under normal versus shear loading that are sampled by the underlying sensors
(Fig.~1d). Sparse tri-axis sampling from a local $2\times2$ unit, a $1\times4$
strip and the full 44-node array is combined with calibrated spatial
reconstruction to move from local decoding to whole-hand mapping (Fig.~1e).
The integrated system comprises 44 MLX90393 nodes with
\dataneeded{number of effective/active nodes} active channels covering
\dataneeded{palmar coverage area, cm$^2$}, with a total thickness of
\dataneeded{system thickness, mm} and mass of \dataneeded{system mass, g}.
It streams a local $2\times2$ frame at \dataneeded{2x2 frame rate, Hz} and a
full 44-node frame at \dataneeded{full-hand frame rate, Hz}, with a dropped-frame
rate of \dataneeded{dropped-frame \% over 10 min} over 10~min of continuous
acquisition and a power draw of \dataneeded{system power, mW}.

% ----------------------------------------------------------------------
\subsection*{Geometry-dependent amplification and physical characterization of magnetic cilia}
% -> Fig. 2
We fabricated NdFeB--elastomer cilia by mixing, degassing, mould filling,
magnetic pre-alignment, curing and demoulding, followed by final magnetization
(Fig.~2a,b; Methods). Cilia were produced across a geometry matrix of lengths
$L=3,5,8$~mm and diameters $d=0.5,0.8$~mm at a pitch of
\dataneeded{measured pitch $s$, mm} (Fig.~2c). For each geometry we measured
\dataneeded{$n$ cilia measured per geometry, target $\geq 30$ from $\geq 3$
patches} cilia and report mean\,$\pm$\,s.d.\ of height, diameter and pitch.
Mechanical characterization of non-magnetic and $x$/$y$/$z$-magnetized
specimens ($n=$~\dataneeded{independent specimens per group, target 5 (>=)})
gives a modulus of \dataneeded{modulus, kPa} and loading--unloading hysteresis
of \dataneeded{mechanical hysteresis, \%} (Fig.~2d).

To test whether the cilia geometry enhances weak-stimulus response, we compared
each of the six geometries against matched controls---a flat magnetic film,
short magnetic pillars, non-magnetic cilia and a bare MLX90393---under
\dataneeded{normal-force levels, e.g.\ 10/20/50/100~mN} normal and
\dataneeded{shear-force levels} shear loads
($n=$~\dataneeded{independent patches, target 3} patches,
\dataneeded{repeats per condition} repeats). We define the force sensitivity
$S=\mathrm{d}|\Delta\mathbf{B}|/\mathrm{d}F$, the cilia gain
$G_{\mathrm{cilia}}=S_{\mathrm{cilia}}/S_{\mathrm{flat}}$ and the limit of
detection $\mathrm{LOD}=3\sigma_B/S$. The optimal geometry
\dataneeded{best $L\times d$} reaches a sensitivity of \dataneeded{$S$, units}
and a gain of \dataneeded{$G_{\mathrm{cilia}}$ with 95\% CI} over the matched
flat film, with an LOD of \dataneeded{LOD, mN} (Fig.~2e). Dynamic
characterization gives a rise time of \dataneeded{rise time, ms}, a recovery
time of \dataneeded{recovery time, ms} and a flat frequency response up to
\dataneeded{bandwidth, Hz}; durability over \dataneeded{cycles, target
$\geq$10{,}000} cycles and baseline drift of \dataneeded{drift over $\geq$30~min}
confirm stability (Fig.~2f).

\noindent\textit{Claim gate.} The term ``cilia-amplified'' is used only if the
cilia group significantly exceeds the flat control with a defined gain across
$\geq 3$ independent devices and statistical significance, after controlling for
total magnetic material and no-load field; otherwise we use ``cilia-enhanced''
or ``magnetic cilia tactile skin'' (see notes/claims\_table.md).

% ----------------------------------------------------------------------
\subsection*{Local tri-axis sensing and unseen-position inverse decoding}
% -> Fig. 3a–e
A local $2\times2$ unit of four MLX90393 sensors (12 channels: 4 sensors
$\times$ $B_x,B_y,B_z$) at a sensor pitch of \dataneeded{sensor pitch, mm}
(Fig.~3a) was calibrated on a UFactory~xArm with an ATI~Nano17 force/torque
sensor and a non-magnetic probe (Fig.~3b). Tri-axis fingerprints were recorded
over normal forces $F_n=\dataneeded{levels, e.g.\ 10/20/50/100~mN}$, shear
forces $F_s=\dataneeded{levels}$ and directions
$\theta=$~\dataneeded{angles, e.g.\ 0--315 deg}
($n=$~\dataneeded{repeats}), revealing distinct normal/shear responses and a
cross-talk of \dataneeded{normal/shear cross-talk, \%} (Fig.~3c).

For inverse decoding we sampled a \dataneeded{grid, e.g.\ $7\times7=49$}
position grid at \dataneeded{force levels} forces with
\dataneeded{repeats} repeats (\dataneeded{total trials} trials total), splitting
strictly by position (60\% train / 20\% validation / 20\% unseen test) with no
loading-time frame shared across splits (Fig.~3d; Methods). Baseline models
(nearest-neighbour fingerprint, ridge/polynomial regression, random forest and
a compact MLP) were compared in three stages: position only; position + $F_n$;
and position + $F_n$ + $F_x,F_y$. On the unseen-position test set the best model
achieves a localization RMSE of \dataneeded{localization RMSE, mm}
(95th percentile \dataneeded{p95 localization error, mm}), a force error of
\dataneeded{force NRMSE, \%} and a median direction error of
\dataneeded{direction error, deg}, with cross-session performance of
\dataneeded{cross-session error} (Fig.~3e). A localization RMSE below the sensor
pitch is required before ``local spatial resolution'' is claimed.

% ----------------------------------------------------------------------
\subsection*{Ultralight bubble-touch perception}
% -> Fig. 3f
As an ultralight-contact demonstration (not a detection-limit measurement), we
released \dataneeded{number of bubbles, target $\geq 20$} soap bubbles across
\dataneeded{bubble-size ranges, target 3} size ranges onto the local unit,
together with \dataneeded{sham trials, target $\geq 2$} no-contact sham trials,
under synchronized video (Fig.~3f). Approach, contact, deformation and release
produced a peak $|\Delta\mathbf{B}|$ of \dataneeded{peak $|\Delta B|$, units}
at an SNR of \dataneeded{SNR, dB}, and predicted contact positions matched the
video frames. A bubble of estimated mass \dataneeded{estimated bubble mass, mg}
was reliably detected and localized; we do not interpret this as a calibrated
force detection limit.

% ----------------------------------------------------------------------
\subsection*{Full-hand integration, artifact suppression and spatial tactile mapping}
% -> Fig. 4
The local capability was scaled to a 44-node hand-shaped FPCB spanning
fingertip, finger pad, palm and thumb regions with per-node coordinates and
orientations (Fig.~4a,b; Extended Data). Because whole-hand maps can be
corrupted by posture and field artifacts, we recorded no-contact motions---open
hand, individual finger flexion, fist, wrist pitch/roll, whole-hand rotation and
cable movement ($n=$~\dataneeded{repeats per motion, target 10})---alongside
standard contacts (palm 50~mN, fingertip 30~mN, thumb 30~mN). We quantify a
contact-to-motion ratio $\mathrm{CMR}=20\log_{10}(A_{\mathrm{contact}}/
A_{\mathrm{motion}})$ of \dataneeded{CMR before correction, dB}, improved to
\dataneeded{CMR after correction, dB} after reference-sensor compensation
$\Delta\mathbf{B}_{i,\mathrm{corr}}=\Delta\mathbf{B}_i-\alpha_i
\Delta\mathbf{B}_{\mathrm{ref}}$, reducing the false-positive rate to
\dataneeded{false-positive rate, \%} (Fig.~4c).

On a fixed hand model we validated \dataneeded{ground-truth positions, target 30}
known positions at \dataneeded{force levels} forces
(\dataneeded{total trials} trials), obtaining a localization error of
\dataneeded{full-hand localization error, mm}, a region confusion accuracy of
\dataneeded{region accuracy, \%} and a false-hotspot ratio of
\dataneeded{false-hotspot ratio} (Fig.~4d). Whole-hand response maps for palm
press, fingertip pinch, cylinder and ball grasp, lateral shear and gentle
contact ($n=$~\dataneeded{repeats per scenario, target 10}) are shown on a common
color scale (Fig.~4e); force-calibrated $F_n,F_x,F_y$ maps are reported only
after per-node calibration, otherwise $|\Delta\mathbf{B}|$ and axis-resolved
maps are shown. Conformability under bending and wrapping, and grasp-state
discrimination across fist, cylinder and sphere, are characterized with a
sensitivity change of \dataneeded{sensitivity change after bending, \%} and a
sensor failure rate of \dataneeded{failure rate after bending cycles, \%}
(Fig.~4f). Absent feedback experiments, we describe this as grasp-state
discrimination and distributed contact readout rather than closed-loop control.

% ----------------------------------------------------------------------
\subsection*{Perception of weak airflow and physiological pulse micro-vibrations}
% -> Fig. 5
Finally, we tested whether one platform senses both external environmental
perturbation and internal physiological micro-vibration. For airflow, we
calibrated the response against a reference anemometer over wind speeds
\dataneeded{wind speeds, e.g.\ 0--5~m\,s$^{-1}$}
(\dataneeded{runs per speed, target 5} runs, randomized order), obtaining a
$\Delta B$--$v$ relationship of \dataneeded{fit form and $R^2$}, a response time
of \dataneeded{airflow response time, s} and up/down hysteresis of
\dataneeded{airflow hysteresis, \%} (Fig.~5a). At a fixed speed of
\dataneeded{fixed speed, m\,s$^{-1}$} across hand poses (palm flat, fingers
flexed, side, back, partial fist; $n=$~\dataneeded{repeats}) we mapped the
spatial distribution of airflow-induced tactile responses---not a resolved flow
field---(Fig.~5b), and tracked dynamic step, staircase and random airflow with a
correlation of \dataneeded{airflow tracking correlation} (Fig.~5c).

For pulse sensing (under ethics approval and informed consent; Methods), a
palmar/wrist patch recorded radial-artery micro-vibrations synchronized with
ECG in \dataneeded{number of subjects, target 10--20} subjects across
\dataneeded{sessions per subject} sessions with re-attachment (Fig.~5d).
Beat segmentation gives a heart-rate MAE of \dataneeded{HR MAE, bpm} against
ECG, an ECG-to-pulse delay of \dataneeded{pulse transit delay, ms}, a waveform
SNR of \dataneeded{pulse SNR, dB} and beat-to-beat CV of \dataneeded{beat CV, \%}
(Fig.~5e). Multi-position amplitude maps (centre, radial, ulnar, proximal,
distal offsets) show the spatial dependence of the pulse signal (Fig.~5f). The
dicrotic notch is reported only where it is stable across beats and subjects.
